% appendix/E-datasheet.tex — Datasheet for the corpus (Gebru et al. 2021 style)
\chapter{Datasheet สำหรับชุดข้อมูล}
\label{app:datasheet}

ภาคผนวกนี้จัดทำในรูปแบบ \emph{Datasheet for Datasets} ตาม Gebru et~al.\
(Communications of the ACM, 2021) เพื่ออธิบาย corpus ที่ใช้ในการศึกษานี้
ในระดับรายละเอียดที่งานวิจัยอื่นสามารถนำไปประเมินความเหมาะสมต่อภาระงาน
ของตนได้

\section{แรงจูงใจ (Motivation)}

\paragraph{Q1. ใครสร้างชุดข้อมูลนี้และเพื่อจุดประสงค์ใด}
ชุดข้อมูลถูกสร้างโดยคณะผู้วิจัยปริญญาตรี ภาควิชาสถิติประยุกต์และการวิเคราะห์
ข้อมูล KMITL เพื่อการศึกษาเปรียบเทียบประสิทธิภาพของแบบจำลองทรานส์ฟอร์เมอร์
ภาษาไทยในงานพยากรณ์ความไวรัล

\paragraph{Q2. ใครเป็นผู้สนับสนุนการสร้างชุดข้อมูล}
งานวิจัยใช้ public API ของ YouTube (free tier 10{,}000 units/day) ไม่มี
แหล่งทุนภายนอก ค่าใช้จ่าย cloud GPU (\$0.49 USD) คณะผู้วิจัยรับผิดชอบเอง

\paragraph{Q3. ตอบช่องว่างใดในวรรณกรรม}
(1) ไม่มีงานก่อนหน้าที่เปรียบเทียบ Thai transformer encoders บนงาน virality
prediction บนชุดข้อมูลเดียวกัน (2) ไม่มี Thai YouTube corpus สาธารณะที่
ออกแบบเพื่อการศึกษาความไวรัลแบบ \emph{intra-channel}

\section{องค์ประกอบ (Composition)}

\paragraph{Q4. ตัวอย่างเป็นอะไร}
หนึ่งตัวอย่าง = หนึ่งวิดีโอ YouTube ภาษาไทย ระดับ public visibility
ประกอบด้วยข้อความ (title, description, tags), meta-data ของช่อง,
ตัวชี้วัด engagement (views, likes, comments) ณ snapshot date

\paragraph{Q5. มีกี่ตัวอย่าง}
23{,}431 วิดีโอ จาก 82 ช่อง เผยแพร่ระหว่าง พ.ศ.\ 2562--2569 (10 ปี)
ขนาดทางสถิติ: ตัวอย่างต่อช่องเฉลี่ย = 285 (ต่ำสุด 12, สูงสุด 1{,}847)

\paragraph{Q6. แต่ละตัวอย่างประกอบด้วยข้อมูลใด}
ดูตารางที่ \ref{tab:fields} ในบทที่ \ref{ch:method} (15 ฟิลด์หลัก)
หลังการแปลงเป็นฟีเจอร์ ผลคือ 27 structural + 6 sentiment + 768 embedding
+ 20{,}000 TF-IDF dimensions (ใช้แยกตามชนิดของแบบจำลอง)

\paragraph{Q7. มี label หรือ target หรือไม่}
มี label เป็นทวิภาค $y_i \in \{0, 1\}$ คำนวณจากดัชนี $\mathit{VI}$ ในบทที่
\ref{ch:method} \S\ref{sec:method-vi} เกณฑ์ = top-decile per channel

\paragraph{Q8. มีข้อมูลที่หายหรือไม่}
มี \texttt{description = ""} ใน $\sim$8\% ของวิดีโอ \texttt{tags = []} ใน
$\sim$15\% มีการ impute เป็นสตริงว่างก่อนสร้างฟีเจอร์ ส่วน \texttt{view\_count = NULL}
ถูกตัดออกในระยะการ clean

\paragraph{Q9. ความสัมพันธ์ระหว่างตัวอย่างมีหรือไม่}
\textbf{มี} วิดีโอที่อยู่ในช่องเดียวกันมีความสัมพันธ์เชิงโครงสร้าง (รูปแบบภาษา
ของผู้สร้าง audience overlap, recommendation algorithm bias) คณะผู้วิจัย
จัดการความสัมพันธ์นี้ผ่าน channel-grouped split (ดู \S\ref{sec:method-split})

\paragraph{Q10. มี split ที่แนะนำหรือไม่}
ใช่ — train (6{,}964) / val (3{,}143) / test (3{,}510) แบบ channel-grouped
+ time-aware (\S\ref{sec:method-split}) split files อยู่ใน
\texttt{data/processed/dataset\_with\_labels.parquet} โดยฟิลด์ \texttt{split}

\paragraph{Q11. มี noise / errors ที่ทราบหรือไม่}
(1) Unicode mojibake ในชื่อวิดีโอบางส่วนของช่องเก่า (พบ $<$1\%) (2) จำนวน
\texttt{like\_count} อาจไม่ตรงกับช่วงเวลาของ \texttt{view\_count} เพราะ API
ส่ง snapshot ณ เวลาที่ query แต่ไม่มีการ versioning

\paragraph{Q12. ชุดข้อมูล self-contained หรือพึ่ง resource ภายนอก}
ครบสมบูรณ์ในไฟล์ \texttt{parquet} เดียว ไม่ต้องการ resource ภายนอก ยกเว้น
ฟีเจอร์ที่สกัดหลังจากนั้น (sentiment cache + title embeddings) ซึ่งสามารถ
recompute ได้จาก raw text

\paragraph{Q13. มีข้อมูลที่อาจถือว่าเป็น secret หรือไม่}
ไม่มี. ทุกฟิลด์ถูกเปิดเผยต่อสาธารณะผ่าน YouTube web UI

\paragraph{Q14. มีข้อมูลที่อาจสร้างความรู้สึกไม่สบายใจให้ผู้อ่านหรือไม่}
ในชื่อวิดีโอบางส่วนของหมวด Entertainment มีคำหยาบหรือคำที่กระตุ้นอารมณ์
ในเชิงลบ คณะผู้วิจัยไม่ได้กรองออก เพราะคำเหล่านี้เป็นส่วนหนึ่งของ phenomenon
ที่ศึกษา

\section{การเก็บรวบรวม (Collection)}

\paragraph{Q15. กระบวนการเก็บข้อมูล}
ใช้ \texttt{src/data\_collection/youtube\_collector.py} เรียก
YouTube Data API v3 endpoint \texttt{videos.list}, \texttt{channels.list},
\texttt{search.list} โดย \texttt{maxResults=50} ต่อ request

\paragraph{Q16. ใครเก็บข้อมูล}
คณะผู้วิจัยเอง ระหว่าง พ.ค.\ 2569

\paragraph{Q17. การจ่ายเงินตอบแทนผู้เข้าร่วม}
ไม่เกี่ยวข้อง (ข้อมูลสาธารณะ ไม่มี human subject)

\paragraph{Q18. ระยะเวลาการเก็บ}
24 ชั่วโมง (single snapshot) ในวันที่ 30 เมษายน 2569

\paragraph{Q19. การพิจารณาทางจริยธรรม / IRB}
ไม่เข้าข่าย IRB เพราะ (1) ไม่มี human subject (2) ใช้ข้อมูลสาธารณะ (3) ไม่มี
PII คณะผู้วิจัยปฏิบัติตาม YouTube API Services Terms of Service ทุกประการ

\section{การ Preprocessing / Labeling / Cleaning}

\paragraph{Q20. การ preprocess ที่ทำ}
ดู \S\ref{sec:method-labels} (clean) และ \S\ref{sec:method-features}
(feature extraction) ทุกขั้นตอนใน \texttt{src/data\_processing/} และ
\texttt{src/features/}

\paragraph{Q21. ข้อมูลดิบยังถูกเก็บไว้หรือไม่}
\textbf{ใช่} \texttt{data/raw/youtube\_thai\_videos.parquet} ถูกเก็บไว้
ในเครื่องของคณะผู้วิจัย และในที่เก็บข้อมูลส่วนตัว (private dataset
repository) ของคณะผู้วิจัยและอาจารย์ที่ปรึกษา

\paragraph{Q22. software ที่ใช้ preprocess}
Python 3.11, pandas 2.x, pyarrow 15.x, PyThaiNLP 5.x ทั้งหมดเป็น
open-source

\section{การใช้งาน (Uses)}

\paragraph{Q23. มีงานวิจัยอื่นที่ใช้ชุดข้อมูลนี้แล้วหรือยัง}
ยังไม่มี ปัญหาพิเศษนี้เป็นการใช้งานครั้งแรก

\paragraph{Q24. ชุดข้อมูลถูกใช้ในงานประเภทใดได้บ้าง}
นอกจาก virality prediction แล้ว สามารถใช้ในงาน (1) Thai title
generation/rewriting (2) channel-level engagement modelling (3) Thai
sentiment-engagement correlation studies

\paragraph{Q25. มีการใช้งานที่ \emph{ไม่ควร} ทำหรือไม่}
\textbf{ไม่ควร} ใช้เพื่อ
(1) เปิดเผยข้อมูลรายบุคคลของผู้สร้างเนื้อหา (channel\_id เปิดเผย แต่ไม่ควร
นำไปจับคู่กับข้อมูลส่วนตัว) (2) สร้างเครื่องมือ \emph{predict-and-attack}
ช่อง competitor (3) optimisation ที่ละเมิด YouTube TOS (เช่น
view-count manipulation)

\section{การจัดเก็บและการแจกจ่าย (Distribution)}

\paragraph{Q26. ชุดข้อมูลจะแจกจ่ายให้บุคคลที่สามหรือไม่}
\textbf{เฉพาะคำขอเชิงวิชาการ} เพื่อปฏิบัติตาม YouTube TOS Section II.B.3
ที่ห้าม redistributable ของ Content เชิงพาณิชย์ ผู้ขอใช้ต้องลงนามใน data
use agreement และมี institutional affiliation

\paragraph{Q27. แจกจ่ายผ่านช่องทางใด}
ผ่านที่เก็บข้อมูลส่วนตัวของคณะผู้วิจัยและอาจารย์ที่ปรึกษา โดยให้สิทธิ์
เป็นรายบุคคลตามเงื่อนไขข้อ Q26

\paragraph{Q28. มี IP หรือสัญญาเงื่อนไขการใช้หรือไม่}
ผู้ขอใช้ต้องอ่านและยอมรับ \emph{YouTube API Services Terms of Service} ก่อน
และตกลงไม่ redistribute ภายใต้เงื่อนไข non-commercial / academic-only

\section{การบำรุงรักษา (Maintenance)}

\paragraph{Q29. ใครรับผิดชอบดูแล / อัปเดต}
ระหว่างที่ปัญหาพิเศษอยู่ในการประเมิน — คณะผู้วิจัย หลังจากนั้น — อาจารย์
ที่ปรึกษา (\advisor)

\paragraph{Q30. มีกำหนดอัปเดตเวอร์ชันใหม่หรือไม่}
ไม่มี ชุดข้อมูลถูก freeze ไว้ตามวันที่ snapshot สำหรับการ reproduce
สถานะปัจจุบันอย่างถาวร

\paragraph{Q31. ติดต่อผู้ดูแลได้ที่ใด}
ผ่านอาจารย์ที่ปรึกษา (\advisor) หรือคณะผู้วิจัยทั้งสาม
(\authorone, \authortwo, \authorthree) ผ่านช่องทางอีเมลของ
\school\ \university
